\section{Pattern Positions}
\label{sec:cses-pattern-positions}

\problemstatement{Given a text $t$ and $k$ patterns, output for each
pattern the starting position of its \emph{first} occurrence in $t$
(1-indexed), or $-1$ if it does not occur.}

\begin{patternbox}
Same suffix automaton, one more annotation: store, for each state, the
ending position of the \textbf{first} occurrence of its substrings
(\textit{firstpos}). A pattern's first-occurrence start is then
$\textit{firstpos}-\lvert p\rvert+1$.
\end{patternbox}

\subsection*{Tracking the first endpos}

When the automaton is extended with the character at index $i$ (1-indexed),
the newly created state's longest substring first \emph{ends} at position
$i$, so set its $\textit{firstpos}=i$. A cloned state inherits the
$\textit{firstpos}$ of the state it was split from, because the clone
represents shorter substrings that first appear no later. After building,
a pattern $p$ that walks to state $v$ has its earliest occurrence ending at
$\textit{firstpos}[v]$, hence starting at
$\textit{firstpos}[v]-\lvert p\rvert+1$.

\begin{ideabox}[The First Endpoint Locates the First Occurrence]
Every substring in a state first occurs at the same earliest ending
position (that is what an endpos class guarantees for its shortest common
context). Recording that position at creation time -- and copying it to
clones -- gives $O(1)$ first-occurrence lookups after the walk.
\end{ideabox}

\subsection*{Algorithm}

\begin{enumerate}
  \item Build the SAM; when creating the state at step $i$, set
        $\textit{firstpos}=i$; for a clone, copy the original's
        $\textit{firstpos}$.
  \item For each pattern, walk to its state $v$; if the walk fails, output
        $-1$; else output $\textit{firstpos}[v]-\lvert p\rvert+1$.
\end{enumerate}

\subsection*{Dry run}

\dryrun{$t=$``aybabtu''; pattern ``ab''.}

\begin{itemize}
  \item ``ab'' first occurs ending at position $4$ (characters $3$--$4$).
  \item Start $=4-2+1=3$, the 1-indexed first position of ``ab''.
\end{itemize}

\subsection*{C++ implementation}

\begin{lstlisting}
#include <bits/stdc++.h>
using namespace std;

struct SAM {
    struct State { int len, link, firstpos; map<char,int> next; };
    vector<State> st; int last;
    SAM() { st.push_back({0, -1, 0, {}}); last = 0; }
    void extend(char c, int pos) {               // pos = 1-indexed end position
        int cur = st.size();
        st.push_back({st[last].len + 1, -1, pos, {}});
        int p = last;
        while (p != -1 && !st[p].next.count(c)) { st[p].next[c] = cur; p = st[p].link; }
        if (p == -1) st[cur].link = 0;
        else {
            int q = st[p].next[c];
            if (st[p].len + 1 == st[q].len) st[cur].link = q;
            else {
                int clone = st.size();
                st.push_back(st[q]);             // inherits firstpos of q
                st[clone].len = st[p].len + 1;
                while (p != -1 && st[p].next[c] == q) { st[p].next[c] = clone; p = st[p].link; }
                st[q].link = st[cur].link = clone;
            }
        }
        last = cur;
    }
};

int main() {
    string t; int k;
    cin >> t >> k;
    SAM sam;
    for (int i = 0; i < (int)t.size(); i++) sam.extend(t[i], i + 1);

    while (k--) {
        string p; cin >> p;
        int cur = 0; bool ok = true;
        for (char c : p) {
            auto it = sam.st[cur].next.find(c);
            if (it == sam.st[cur].next.end()) { ok = false; break; }
            cur = it->second;
        }
        cout << (ok ? sam.st[cur].firstpos - (int)p.size() + 1 : -1) << "\n";
    }
    return 0;
}
\end{lstlisting}

\begin{complexitybox}
\textbf{Time Complexity.} $O(|t|\log\Sigma)$ build, $O(\sum|p_i|
\log\Sigma)$ queries. \textbf{Space Complexity.} $O(|t|)$.
\end{complexitybox}

\begin{takeawaybox}
Recording each state's first ending position turns the SAM into a
first-occurrence locator. Combined with endpos counts (previous section),
one automaton answers presence, frequency, and position -- the full
substring query trio.
\end{takeawaybox}
